\enteteExam{NFE106 -- Administration de bases de données}%
{Examen 22 juin 2022 -- session 1 -- 3h}%
{Examen sur table -- Seuls documents autorisés\,: 10 pages maximum de notes manuscrites\,; 
calculettes autorisées}%


\section{Indexation et tri (6 points)}

Voici le contenu d'un fichier \texttt{animaux}\,:

\begin{verbatim}
(jaguar, 17), (chameau, 22), (gnou, 1), (girafe, 13), (chat, 3),
(lion, 8), (dauphin, 40), (zèbre, 11), (hamster, 6), (hyène,9),
(licorne, 2), (babouin, 12), (piranha,55), (saumon, 82), (bar,
44), (anguille, 43), (gorille,98), (tigre,76), (requin, 56), (escargot,45).
\end{verbatim}

On suppose qu'on peut placer  2 enregistrements par bloc. 

\begin{enumerate}
  \item On veut construire un index non-dense sur le premier attribut.
    Que faut-il faire au préalable\,? Donnez les différents niveeaux de l'index.
 
  \item Construire un arbre B sur le second attribut, 
     en supposant 2 enregistrements et trois pointeurs par bloc \emph{au maximum}.

  \item On dispose des deux index ci-dessus. Quel est le 
     nombre d'entrées/sorties \emph{dans le pire des cas}
     pour la recherche des animaux dont le nom commence par un 'g',
     et pour la recherche des animaux dont l'identifiant est compris
    entre 40 et 50 (prendre en compte les accès à l'index \emph{et}
    au fichier).

  \item On veut trier ce fichier (tel qu'il est donné dans l'énoncé)
     avec seulement 3 blocs, toujours
    en supposant qu'on peut placer deux enregistrements par bloc.
    Décrivez le déroulement de l'algorithme de tri-fusion, 
   et donnez le nombre total d'entrées/sorties, sans compter
    l'écriture finale du fichier trié. 
\end{enumerate}

\begin{solution}
\begin{enumerate}
  \item Il faut d'abord trier le fichier:
\begin{verbatim}
(anguille, 43),  (bar, 44), (babouin, 12), (chameau, 22),  
(chat, 3), (dauphin, 40),  (escargot,44),
(girafe, 13), (gnou, 1),  (gorille,98), 
(hamster, 6), (hyène,9), (jaguar, 17),  
(licorne, 2),  (lion, 8),   (piranha,55), (saumon, 82),  (requin, 56),
 (tigre,76), (zèbre, 11),
\end{verbatim}
 Ensuite on groupe par 2 et on crée le premier niveau d'index:
\begin{verbatim}
anguille -- babouin -- chat -- escargot -- gnou -- hasmter --
jaguar -- lion -- saumon -- tigre
\end{verbatim}

Encore une fois:
\begin{verbatim}
anguille --  chat --  gnou -- jaguar -- saumon 
\end{verbatim}

Puis
\begin{verbatim}
anguille --   gnou -- saumon 
\end{verbatim}

Et finalement il reste \texttt{(anguille, saumon)} à la racine, soit 4 niveaux.

  \item Standard.

  \item La première recherche doit ramener trois enregistrements. 
  Dans le cas de l'index non-dense, le fichier est trié.
     On peut donc parcourir séquentiellement à partir du
  premier 'g'. On lit donc 4 blocs dans l'index pour arriver aux feuilles, 
  puis 2 blocs en parcours séquentiel.
  
  Les enregistrements ne sont pas ordonnés pour l'arbre B. Il faut donc, après le
  parcours d'index, parcourir les feuilles pour les valeurs de clé entre 40 et 50,
  avec une lecture de bloc (au pire) pour chacune.

  \item On utilise 2 blocs pour l'entrée, un pour la sortie. Dans
   les 2 blocs on place 4 enregistrements. Il y a 20 animaux
   donc (1) 5 fragments initiaux de 2 blocs, puis (2)
    2 fragments de 4 blocs, e
    un  fragment de 2 bloc, puis (3) un fragment de 8 blocs et un  de deux.
   Et on termine par une lecture. Donc coût =$3 \times 2 \times 10 + 10 = 70$.
\end{enumerate}
\end{solution}

\section{Jointures et optimisation (8 points)}

On considère trois relations 
$R(\textbf{a},b,d)$, $S(\textbf{c},d,e)$ et $T(\textbf{e},f)$ dont les clés
primaires sont respectivement $a$, $c$ et $e$.
$R$ contient 200~000 enregistrements, $S$  20 enregistrements et $T$ 500 enregistrements. Des index
sont créés sur les clés primaires, et on suppose que
pour chaque relation, y compris celles qui sont calculées par une jointure,
on stocke 10 enregistrements par bloc.

On suppose que tous les index sont toujours en mémoire principale
(donc pas d'entrée/sortie pour les accès aux index).
On dispose de 20 blocs en mémoire pour traiter les jointures.
Les jointures se font sur les attributs de même nom.

\begin{enumerate}
  \item Quelle est le nombre maximal
    de enregistrements de $S \Join T$ (indiquez également la condition 
     pour que ce nombre maximal soit atteint)\,? Quelle est la clé de la
    relation obtenue\,?

  \item Mêmes questions pour $R \Join S \Join T$.

  \item Décrire le fonctionnement de l'algorithme par boucles 
   imbriquées pour calculer $S \Join T$, en exploitant au mieux la
   mémoire disponible. 
   Indiquez le coût en entrées/sorties pour cet algorithme?

  \item Même question l'algorithme par boucles 
   imbriquées \emph{indexées}.
   
  \item Décrire un plan d'exécution  
    pour calculer $R \Join S \Join T$\ et 
    évaluer son coût (nombre de blocs lus).
    
    \item  Une application a inséré des enregistrements dans $S$
qui contient maintenant 5~000 enregistrements. Le
SGBD gère un histogramme qui indique que la sélectivité
de l'attribut $d$ est 5 (autrement dit, une sélection 
sur $R$ ou $S$ pour une valeur de $d$ ramène 5\% des enregistrements).
Quelle taille peut-on estimer pour $R \Join S$\,?

  \item (\textbf{Question bonus, pour 2 points}) Je n'ai toujours que 20 blocs en mémoire.
  Décrire l'algorithme de jointure par hachage pour $R \Join S$ 
  et évaluer son coût.
\end{enumerate}

\begin{solution}
 \begin{enumerate}
   \item Il y a au plus 20 enregistrements, et il faut pour l'atteindre
       que la clé étrangère soit \texttt{not null},
       et le respect de l'intégrité référentielle de $S$ vers $T$.

  \item Dans le pire des cas il y a $200~000 \times 20$ enregistrements
        (s'il y a une seule valeur pour $d$).

   \item On met $S$ en mémoire (donc 2 blocs) et on l'organise
   comme une table de hachage sur l'attribut $S.e$. Puis on lit $T$ séquentiellement
   (50 blocs),
   en cherchant pour chaque enregistrement $(e_i,f_i)$ de $T$ le ou les
   enregistrements correspondant dans la table de hachage.
   
   Coût\,,: 52 blocs.
   
   \item On parcourt $S$ séquentiellement (2 blocs), pour chaque enregistrement
   $(c, d,e)$ (il y en a 20), on utilise l'index pour accéder à $T$. Pour chaque entrée
   trouvée dans l'index il faut lire un bloc (au pire) sur le disque.
   
       Coût\,,: 2 + 20 blocs. 

   \item On constate qu'il n'existe pas d'index disponible 
      pour la jointure avec $R$, car la clé primaire de $R$ ne se trouve
      en tant que clé étrangère ni dans $S$ ni dans $T$.
      
      L'algorithme\,: on évalue  $S \Join T$, comme indiqué précemment. 
      Il faut ensuite choisir un algorithme de jointure sans index.
      Etant donnée la petite taille de $S \Join T$, on choisit
     la jointure par boucles imbriquées.
     On construit une table de hachage en mémoire sur le résultat de $S \Join T$
     (2 blocs au maximum, clé de hachage\,: $d$)\,; on lit séquentiellement $R$ 
     a joignant chaque enregistrement avec la table de hachage. 
     
     La lecture de $R$ est prédominante\,: 20 000 blocs à lire séquentiellement.
     
     Toutes les autres possibilités sont moins bonnes, à cause de la taille
     du résultat de la jointure entre $R$ et une autre table. Penser par exemple
     au coût qui résulterait d'un tri préalable de $R$ sur l'attribut $d$ pour
     une jointure par tri-fusion.
     
   \item Pour chaque tuple de $R$, on trouve $5000 \times 0,05 = 250$
    enregistrements dans $S$: 50 M de enregistrements au final.

   \item Jointure par hachage\,:je hache $S$ en 25 fragments de 
   200 enresgistremets chacune (20 blocs en mémoire). Ensuite je hache $R$ en 25
   fragments également (leur taille m'importe peu). Puis je fais la jointure
   entre chaque paire de fragmentss. Coût =
   $3 \times (|S| + |R|) = 3 \times (20\,000 + 500) = 61\,500$.
 
 \end{enumerate}
\end{solution}


%%%%%%%%%%%%%%%%%%%%%%%%

\section{Concurrence (6 points)}

On considère le système d'information
d'un institut de sondage, avec les tables relationnelles 
suivantes (les attributs ou combinaisons d'attributs qui forment 
une clé unique sont \underline{soulignés}).

\begin{itemize}
  \item  Personne (\textbf{numPers}, nom, sexe, numCat)
  \item Question (\textbf{numQ}, description)
  \item  Avis (\textbf{numQ, numPers}, reponse)
\end{itemize}


L'exécution suivante est reçue par le système de l'institut de sondage :
$$H : r_1[x] r_2[y] w_1[x] r_3[y]  r_2[x] w_3[y] r_2[z] C_1 r_3[z] w_2[z] C_2 w_3[z] C_3$$


\begin{enumerate}
  \item Parmi les programmes qui s'exécutent dans le système, 
il y a \texttt{ModifierAvis(nomPers, numQuestion, avis\_donné,
nouvel\_avis)},  qui modifie 
la réponse \textit{réponse\_donnée} par
la personne \texttt{nomPers} à  la question \texttt{numQuestion} 
en \textit{nouvelle\_réponse}. 

Si
les enregistrements de $H$ sont des nuplets des relations de la base 
de données, montrez (en justifiant votre réponse)
quelles transactions de $H$ pourraient provenir de \textit{ModifierAvis}.



  \item Vérifiez si $H$ est sérialisable en identifiant 
les conflits et en construisant le graphe de sérialisation.

  \item Montrer qu'il existe des lectures sales, 
    et expliquez les conséquences possibles.
 
  \item  Quelle est l'exécution obtenue par verrouillage 
 à deux phases à partir de $H$? 

  \item (\textbf{Question bonus, pour 2 points}). Quelle est l'exécution obtenue avec
    l'algorithme de concurrence par versionnement\,?
         On suppose que toutes les transactions débutent
       au même moment.

\end{enumerate}


\begin{solution}

\begin{enumerate}
  \item $T_1 = r_1[x] w_1[x]$, $T_2 = r_2[y] r_2[x] r_2[z] w_2[z]$,
           $T_3 = r_3[y] w_3[y] r_3[z] w_3[z]$  
           
           \texttt{ModifierAvis} doit lire un avis et 
           le modifier. On peut donc considérer qu'il s'agit de $T_1$.

   \item $r_2[x]$ est une lecture sale (car précédée de $w_1[x]$). 
   Si $T_1$ s'interrompt avant le \texttt{commit}, la valeur
   lue par $r_2[x]$ n'existera plus et $T_2$  risque de 
   valider un état incohérent.

   \item $r_1[x] r_2[y] w_1[x] r_3[y] (T_2\, attente sur r_2[x]) (T_3\ attente sur w_3[y])
      C_1 r_2[x] r_2[z] w_2[z] C_2 w_3[y] r_3[z] w_3[z] C_3$

   \item Avec l'algo de versionnement, on rejette $w_3[z]$ 
    car l'écriture $w_2[z]$ est intervenue entre-temps.
\end{enumerate}

\end{solution}

